\clearpage
\section{Two-VM remote reads and writes (run 9)}\label{sec:rwc}
\textbf{\RWCstate} \RWCFINDING{parity}

\textbf{Topology and parity.} Selection \texttt{remote\_rw\_compact} uses one coordinator and one remote shard-WAL VM, both \texttt{z4m-highmem-16}, in \texttt{us-central1-a} with compact placement. Native runs in process on the coordinator SSD; owner TCP/verbs use the remote owner SSD. Workloads match size, caller count and repetition; read arms also match keys/order, geometry, distribution, operation/warmup counts, mix, read tier and zero coordinator-cache bytes within each temperature. Network/security, index/history, routing and host SSD remain different. Run-9 TCP used mutual TLS and verbs trusted/no TLS, so their ratios include the security difference; run 10 adds trusted TCP rather than treating run 9 as like-for-like transport evidence.

\textbf{Lanes and acknowledgement.} Native has zero remote-read lanes (in-process). Each owner uses one lane at c1/c5 and four at c16. Thus c5 owner reads are lane-limited, \emph{not a transport limit or a basis to rank transports}. Native page cache is not durable; Native per-put file/directory fsync, grouped Native barrier and owner covering file fdatasync remain distinct. Similar durable intent does not establish equal crash recovery; RF1 is not failover.

Cold eligibility requires a qualified receipt and at least 0.9 device reads per successful call, not a cache drop alone. This is a case-level proxy, not every-call NAND proof. Native 1 KiB is unqualified at all reader counts and both repetitions; Native 8/64 KiB and every owner cold cell qualify. Warm reads prewarm the serving-node page cache, not the coordinator cache. Cold/warm geometry differs; their contrast is descriptive. No calls, repetitions or fleets are pooled.
\begin{figure}[H]\centering\RWCfigure{parity}
\caption{Run-9 availability: pairs in the first two columns, distinct main-grid cases in the last three. Failed network floors are separate from KV cases. Each owner write has three Native acknowledgement controls, not three equivalent durability contracts.}
\end{figure}
\begin{table}[H]\centering\small\setlength{\tabcolsep}{3pt}
\begin{tabular}{lrrrrr}\toprule
Grid / rep. & Matched pairs & Missing pairs & Failed cases & Unqual. cases & Other cases\\\midrule
Write / r0 & \N{RWC:availability:G14O1:0:matched} & \N{RWC:availability:G14O1:0:missing partner} & \N{RWC:availability:G14O1:0:failed} & \N{RWC:availability:G14O1:0:unqualified cold} & \N{RWC:availability:G14O1:0:not available}\\
Write / r1 & \N{RWC:availability:G14O1:1:matched} & \N{RWC:availability:G14O1:1:missing partner} & \N{RWC:availability:G14O1:1:failed} & \N{RWC:availability:G14O1:1:unqualified cold} & \N{RWC:availability:G14O1:1:not available}\\
Cold read / r0 & \N{RWC:availability:G15C:0:matched} & \N{RWC:availability:G15C:0:missing partner} & \N{RWC:availability:G15C:0:failed} & \N{RWC:availability:G15C:0:unqualified cold} & \N{RWC:availability:G15C:0:not available}\\
Cold read / r1 & \N{RWC:availability:G15C:1:matched} & \N{RWC:availability:G15C:1:missing partner} & \N{RWC:availability:G15C:1:failed} & \N{RWC:availability:G15C:1:unqualified cold} & \N{RWC:availability:G15C:1:not available}\\
Warm read / r0 & \N{RWC:availability:G15M:0:matched} & \N{RWC:availability:G15M:0:missing partner} & \N{RWC:availability:G15M:0:failed} & \N{RWC:availability:G15M:0:unqualified cold} & \N{RWC:availability:G15M:0:not available}\\
Warm read / r1 & \N{RWC:availability:G15M:1:matched} & \N{RWC:availability:G15M:1:missing partner} & \N{RWC:availability:G15M:1:failed} & \N{RWC:availability:G15M:1:unqualified cold} & \N{RWC:availability:G15M:1:not available}\\
\bottomrule\end{tabular}
\caption{Separate controls: Class A Native page cache; Class B Native barriers and owner sync. No cross-class ranking. First two columns count pairs; last three distinct main-grid cases including Native. An owner write has three Native controls, an owner read one. Columns are not a partition. The three failed network-floor cases are separate, not failed KV calls. All sizes and c1/c5/c16 included; pairing does not equalize read lanes or durability. Unmatched dimensions: sync/grouping, RPC/security, index/routing, and host SSD (Native coordinator SSD versus owner SSD). Read lanes: Native 0, owner 1 at c1/c5 and 4 at c16; c5 owner reads are lane-limited, not a transport ranking. Floors are separate populations.}
\end{table}


\clearpage
\subsection{Matched writes: size and acknowledgement}
\begin{figure}[H]\centering\RWCfigure{write-latency}
\caption{Equal writers within each panel; p50 lines and p99 crosses, repetitions separate. Native page cache is non-durable context; other arms wait for per-put, grouped or covering sync. One data SSD per arm, on different hosts.}
\end{figure}
\begin{table}[H]\centering\small\setlength{\tabcolsep}{3pt}
\begin{tabular}{lrrrrr}\toprule
Size & Arm & Rep. & c1 & c5 & c16\\\midrule
1 KiB & Native page cache & r0 & \N{RWC:latency_vs_size:17:p50_us:1|.2f}/\N{RWC:latency_vs_size:17:p99_us:1|.2f} & \N{RWC:latency_vs_size:27:p50_us:1|.2f}/\N{RWC:latency_vs_size:27:p99_us:1|.2f} & \N{RWC:latency_vs_size:37:p50_us:1|.2f}/\N{RWC:latency_vs_size:37:p99_us:1|.2f}\\
1 KiB & Native page cache & r1 & \N{RWC:latency_vs_size:26:p50_us:1|.2f}/\N{RWC:latency_vs_size:26:p99_us:1|.2f} & \N{RWC:latency_vs_size:36:p50_us:1|.2f}/\N{RWC:latency_vs_size:36:p99_us:1|.2f} & \N{RWC:latency_vs_size:46:p50_us:1|.2f}/\N{RWC:latency_vs_size:46:p99_us:1|.2f}\\
1 KiB & Native per-put barrier & r0 & \N{RWC:latency_vs_size:18:p50_us:1|.1f}/\N{RWC:latency_vs_size:18:p99_us:1|.1f} & \N{RWC:latency_vs_size:28:p50_us:1|.1f}/\N{RWC:latency_vs_size:28:p99_us:1|.0f} & \N{RWC:latency_vs_size:38:p50_us:1|.1f}/\N{RWC:latency_vs_size:38:p99_us:1|.0f}\\
1 KiB & Native per-put barrier & r1 & \N{RWC:latency_vs_size:25:p50_us:1|.1f}/\N{RWC:latency_vs_size:25:p99_us:1|.1f} & \N{RWC:latency_vs_size:35:p50_us:1|.1f}/\N{RWC:latency_vs_size:35:p99_us:1|.0f} & \N{RWC:latency_vs_size:45:p50_us:1|.1f}/\N{RWC:latency_vs_size:45:p99_us:1|.0f}\\
1 KiB & Native grouped barrier & r0 & \N{RWC:latency_vs_size:19:p50_us:1|.1f}/\N{RWC:latency_vs_size:19:p99_us:1|.1f} & \N{RWC:latency_vs_size:29:p50_us:1|.1f}/\N{RWC:latency_vs_size:29:p99_us:1|.0f} & \N{RWC:latency_vs_size:39:p50_us:1|.1f}/\N{RWC:latency_vs_size:39:p99_us:1|.0f}\\
1 KiB & Native grouped barrier & r1 & \N{RWC:latency_vs_size:24:p50_us:1|.1f}/\N{RWC:latency_vs_size:24:p99_us:1|.1f} & \N{RWC:latency_vs_size:34:p50_us:1|.1f}/\N{RWC:latency_vs_size:34:p99_us:1|.0f} & \N{RWC:latency_vs_size:44:p50_us:1|.1f}/\N{RWC:latency_vs_size:44:p99_us:1|.0f}\\
1 KiB & Owner TCP & r0 & \N{RWC:latency_vs_size:20:p50_us:1|.0f}/\N{RWC:latency_vs_size:20:p99_us:1|.0f} & \N{RWC:latency_vs_size:30:p50_us:1|.0f}/\N{RWC:latency_vs_size:30:p99_us:1|.0f} & \N{RWC:latency_vs_size:40:p50_us:1|.0f}/\N{RWC:latency_vs_size:40:p99_us:1|.0f}\\
1 KiB & Owner TCP & r1 & \N{RWC:latency_vs_size:23:p50_us:1|.0f}/\N{RWC:latency_vs_size:23:p99_us:1|.0f} & \N{RWC:latency_vs_size:33:p50_us:1|.0f}/\N{RWC:latency_vs_size:33:p99_us:1|.0f} & \N{RWC:latency_vs_size:43:p50_us:1|.0f}/\N{RWC:latency_vs_size:43:p99_us:1|.0f}\\
1 KiB & Owner verbs & r0 & \N{RWC:latency_vs_size:21:p50_us:1|.0f}/\N{RWC:latency_vs_size:21:p99_us:1|.0f} & \N{RWC:latency_vs_size:31:p50_us:1|.0f}/\N{RWC:latency_vs_size:31:p99_us:1|.0f} & \N{RWC:latency_vs_size:41:p50_us:1|.0f}/\N{RWC:latency_vs_size:41:p99_us:1|.0f}\\
1 KiB & Owner verbs & r1 & \N{RWC:latency_vs_size:22:p50_us:1|.0f}/\N{RWC:latency_vs_size:22:p99_us:1|.0f} & \N{RWC:latency_vs_size:32:p50_us:1|.0f}/\N{RWC:latency_vs_size:32:p99_us:1|.0f} & \N{RWC:latency_vs_size:42:p50_us:1|.0f}/\N{RWC:latency_vs_size:42:p99_us:1|.0f}\\
8 KiB & Native page cache & r0 & \N{RWC:latency_vs_size:47:p50_us:1|.2f}/\N{RWC:latency_vs_size:47:p99_us:1|.2f} & \N{RWC:latency_vs_size:57:p50_us:1|.2f}/\N{RWC:latency_vs_size:57:p99_us:1|.2f} & \N{RWC:latency_vs_size:67:p50_us:1|.2f}/\N{RWC:latency_vs_size:67:p99_us:1|.2f}\\
8 KiB & Native page cache & r1 & \N{RWC:latency_vs_size:56:p50_us:1|.2f}/\N{RWC:latency_vs_size:56:p99_us:1|.2f} & \N{RWC:latency_vs_size:66:p50_us:1|.2f}/\N{RWC:latency_vs_size:66:p99_us:1|.2f} & \N{RWC:latency_vs_size:76:p50_us:1|.2f}/\N{RWC:latency_vs_size:76:p99_us:1|.2f}\\
8 KiB & Native per-put barrier & r0 & \N{RWC:latency_vs_size:48:p50_us:1|.1f}/\N{RWC:latency_vs_size:48:p99_us:1|.1f} & \N{RWC:latency_vs_size:58:p50_us:1|.1f}/\N{RWC:latency_vs_size:58:p99_us:1|.0f} & \N{RWC:latency_vs_size:68:p50_us:1|.0f}/\N{RWC:latency_vs_size:68:p99_us:1|.0f}\\
8 KiB & Native per-put barrier & r1 & \N{RWC:latency_vs_size:55:p50_us:1|.1f}/\N{RWC:latency_vs_size:55:p99_us:1|.1f} & \N{RWC:latency_vs_size:65:p50_us:1|.1f}/\N{RWC:latency_vs_size:65:p99_us:1|.0f} & \N{RWC:latency_vs_size:75:p50_us:1|.0f}/\N{RWC:latency_vs_size:75:p99_us:1|.0f}\\
8 KiB & Native grouped barrier & r0 & \N{RWC:latency_vs_size:49:p50_us:1|.1f}/\N{RWC:latency_vs_size:49:p99_us:1|.1f} & \N{RWC:latency_vs_size:59:p50_us:1|.1f}/\N{RWC:latency_vs_size:59:p99_us:1|.0f} & \N{RWC:latency_vs_size:69:p50_us:1|.0f}/\N{RWC:latency_vs_size:69:p99_us:1|.0f}\\
8 KiB & Native grouped barrier & r1 & \N{RWC:latency_vs_size:54:p50_us:1|.1f}/\N{RWC:latency_vs_size:54:p99_us:1|.1f} & \N{RWC:latency_vs_size:64:p50_us:1|.1f}/\N{RWC:latency_vs_size:64:p99_us:1|.0f} & \N{RWC:latency_vs_size:74:p50_us:1|.0f}/\N{RWC:latency_vs_size:74:p99_us:1|.0f}\\
8 KiB & Owner TCP & r0 & \N{RWC:latency_vs_size:50:p50_us:1|.0f}/\N{RWC:latency_vs_size:50:p99_us:1|.0f} & \N{RWC:latency_vs_size:60:p50_us:1|.0f}/\N{RWC:latency_vs_size:60:p99_us:1|.0f} & \N{RWC:latency_vs_size:70:p50_us:1|.0f}/\N{RWC:latency_vs_size:70:p99_us:1|.0f}\\
8 KiB & Owner TCP & r1 & \N{RWC:latency_vs_size:53:p50_us:1|.0f}/\N{RWC:latency_vs_size:53:p99_us:1|.0f} & \N{RWC:latency_vs_size:63:p50_us:1|.0f}/\N{RWC:latency_vs_size:63:p99_us:1|.0f} & \N{RWC:latency_vs_size:73:p50_us:1|.0f}/\N{RWC:latency_vs_size:73:p99_us:1|.0f}\\
8 KiB & Owner verbs & r0 & \N{RWC:latency_vs_size:51:p50_us:1|.0f}/\N{RWC:latency_vs_size:51:p99_us:1|.0f} & \N{RWC:latency_vs_size:61:p50_us:1|.0f}/\N{RWC:latency_vs_size:61:p99_us:1|.0f} & \N{RWC:latency_vs_size:71:p50_us:1|.0f}/\N{RWC:latency_vs_size:71:p99_us:1|.0f}\\
8 KiB & Owner verbs & r1 & \N{RWC:latency_vs_size:52:p50_us:1|.0f}/\N{RWC:latency_vs_size:52:p99_us:1|.0f} & \N{RWC:latency_vs_size:62:p50_us:1|.0f}/\N{RWC:latency_vs_size:62:p99_us:1|.0f} & \N{RWC:latency_vs_size:72:p50_us:1|.0f}/\N{RWC:latency_vs_size:72:p99_us:1|.0f}\\
64 KiB & Native page cache & r0 & \N{RWC:latency_vs_size:77:p50_us:1|.1f}/\N{RWC:latency_vs_size:77:p99_us:1|.1f} & \N{RWC:latency_vs_size:87:p50_us:1|.1f}/\N{RWC:latency_vs_size:87:p99_us:1|.0f} & \N{RWC:latency_vs_size:12:p50_us:1|.0f}/\N{RWC:latency_vs_size:12:p99_us:1|.0f}\\
64 KiB & Native page cache & r1 & \N{RWC:latency_vs_size:86:p50_us:1|.1f}/\N{RWC:latency_vs_size:86:p99_us:1|.0f} & \N{RWC:latency_vs_size:96:p50_us:1|.1f}/\N{RWC:latency_vs_size:96:p99_us:1|.0f} & \N{RWC:latency_vs_size:101:p50_us:1|.0f}/\N{RWC:latency_vs_size:101:p99_us:1|.0f}\\
64 KiB & Native per-put barrier & r0 & \N{RWC:latency_vs_size:78:p50_us:1|.0f}/\N{RWC:latency_vs_size:78:p99_us:1|.0f} & \N{RWC:latency_vs_size:88:p50_us:1|.0f}/\N{RWC:latency_vs_size:88:p99_us:1|.0f} & \N{RWC:latency_vs_size:13:p50_us:1|.0f}/\N{RWC:latency_vs_size:13:p99_us:1|.0f}\\
64 KiB & Native per-put barrier & r1 & \N{RWC:latency_vs_size:85:p50_us:1|.0f}/\N{RWC:latency_vs_size:85:p99_us:1|.0f} & \N{RWC:latency_vs_size:95:p50_us:1|.0f}/\N{RWC:latency_vs_size:95:p99_us:1|.0f} & \N{RWC:latency_vs_size:100:p50_us:1|.0f}/\N{RWC:latency_vs_size:100:p99_us:1|.0f}\\
64 KiB & Native grouped barrier & r0 & \N{RWC:latency_vs_size:79:p50_us:1|.0f}/\N{RWC:latency_vs_size:79:p99_us:1|.0f} & \N{RWC:latency_vs_size:89:p50_us:1|.0f}/\N{RWC:latency_vs_size:89:p99_us:1|.0f} & \N{RWC:latency_vs_size:14:p50_us:1|.0f}/\N{RWC:latency_vs_size:14:p99_us:1|.0f}\\
64 KiB & Native grouped barrier & r1 & \N{RWC:latency_vs_size:84:p50_us:1|.0f}/\N{RWC:latency_vs_size:84:p99_us:1|.0f} & \N{RWC:latency_vs_size:94:p50_us:1|.0f}/\N{RWC:latency_vs_size:94:p99_us:1|.0f} & \N{RWC:latency_vs_size:99:p50_us:1|.0f}/\N{RWC:latency_vs_size:99:p99_us:1|.0f}\\
64 KiB & Owner TCP & r0 & \N{RWC:latency_vs_size:80:p50_us:1|.0f}/\N{RWC:latency_vs_size:80:p99_us:1|.0f} & \N{RWC:latency_vs_size:90:p50_us:1|.0f}/\N{RWC:latency_vs_size:90:p99_us:1|.0f} & \N{RWC:latency_vs_size:15:p50_us:1|.0f}/\N{RWC:latency_vs_size:15:p99_us:1|.0f}\\
64 KiB & Owner TCP & r1 & \N{RWC:latency_vs_size:83:p50_us:1|.0f}/\N{RWC:latency_vs_size:83:p99_us:1|.0f} & \N{RWC:latency_vs_size:93:p50_us:1|.0f}/\N{RWC:latency_vs_size:93:p99_us:1|.0f} & \N{RWC:latency_vs_size:98:p50_us:1|.0f}/\N{RWC:latency_vs_size:98:p99_us:1|.0f}\\
64 KiB & Owner verbs & r0 & \N{RWC:latency_vs_size:81:p50_us:1|.0f}/\N{RWC:latency_vs_size:81:p99_us:1|.0f} & \N{RWC:latency_vs_size:91:p50_us:1|.0f}/\N{RWC:latency_vs_size:91:p99_us:1|.0f} & \N{RWC:latency_vs_size:16:p50_us:1|.0f}/\N{RWC:latency_vs_size:16:p99_us:1|.0f}\\
64 KiB & Owner verbs & r1 & \N{RWC:latency_vs_size:82:p50_us:1|.0f}/\N{RWC:latency_vs_size:82:p99_us:1|.0f} & \N{RWC:latency_vs_size:92:p50_us:1|.0f}/\N{RWC:latency_vs_size:92:p99_us:1|.0f} & \N{RWC:latency_vs_size:97:p50_us:1|.0f}/\N{RWC:latency_vs_size:97:p99_us:1|.0f}\\
\bottomrule\end{tabular}
\caption{Writes; successful API p50/p99 (\us{}), one row per repetition, never pooled. All three caller counts are shown. Values below 10 use two decimals, below 100 one, otherwise whole microseconds. Native page-cache writes are not durable; barrier/grouped/covering sync remain distinct. Unmatched dimensions: sync/grouping, RPC/security, index/routing, and host SSD (Native coordinator SSD versus owner SSD). Read lanes: Native 0, owner 1 at c1/c5 and 4 at c16; c5 owner reads are lane-limited, not a transport ranking. Floors are separate populations.}
\end{table}

\RWCFINDING{writes}

\clearpage
\subsection{Matched reads: device-qualified and prewarmed}
\begin{figure}[H]\centering\RWCfigure{read-latency}
\caption{Matched readers, sizes and repetitions; coordinator cache off. Unqualified Native cold 1 KiB has no latency marker. Native is in process; owners use one lane at c1, four at c16. Appendix~\ref{app:rwc-qualification} retains every verdict and device ratio, including c5.}
\end{figure}
\begin{table}[H]\centering\small\setlength{\tabcolsep}{3pt}
\begin{tabular}{lrrrrr}\toprule
Size & Arm & Rep. & c1 & c5* & c16\\\midrule
1 KiB & Native & r0 & unqualified & unqualified & unqualified\\
1 KiB & Native & r1 & unqualified & unqualified & unqualified\\
1 KiB & Owner TCP & r0 & \N{RWC:latency_vs_size:1:p50_us:1|.0f}/\N{RWC:latency_vs_size:1:p99_us:1|.0f} & \N{RWC:latency_vs_size:106:p50_us:1|.0f}/\N{RWC:latency_vs_size:106:p99_us:1|.0f} & \N{RWC:latency_vs_size:112:p50_us:1|.0f}/\N{RWC:latency_vs_size:112:p99_us:1|.0f}\\
1 KiB & Owner TCP & r1 & \N{RWC:latency_vs_size:103:p50_us:1|.0f}/\N{RWC:latency_vs_size:103:p99_us:1|.0f} & \N{RWC:latency_vs_size:109:p50_us:1|.0f}/\N{RWC:latency_vs_size:109:p99_us:1|.0f} & \N{RWC:latency_vs_size:115:p50_us:1|.0f}/\N{RWC:latency_vs_size:115:p99_us:1|.0f}\\
1 KiB & Owner verbs & r0 & \N{RWC:latency_vs_size:2:p50_us:1|.0f}/\N{RWC:latency_vs_size:2:p99_us:1|.0f} & \N{RWC:latency_vs_size:107:p50_us:1|.0f}/\N{RWC:latency_vs_size:107:p99_us:1|.0f} & \N{RWC:latency_vs_size:113:p50_us:1|.0f}/\N{RWC:latency_vs_size:113:p99_us:1|.0f}\\
1 KiB & Owner verbs & r1 & \N{RWC:latency_vs_size:102:p50_us:1|.0f}/\N{RWC:latency_vs_size:102:p99_us:1|.0f} & \N{RWC:latency_vs_size:108:p50_us:1|.0f}/\N{RWC:latency_vs_size:108:p99_us:1|.0f} & \N{RWC:latency_vs_size:114:p50_us:1|.0f}/\N{RWC:latency_vs_size:114:p99_us:1|.0f}\\
8 KiB & Native & r0 & \N{RWC:latency_vs_size:117:p50_us:1|.2f}/\N{RWC:latency_vs_size:117:p99_us:1|.0f} & \N{RWC:latency_vs_size:123:p50_us:1|.2f}/\N{RWC:latency_vs_size:123:p99_us:1|.0f} & \N{RWC:latency_vs_size:129:p50_us:1|.2f}/\N{RWC:latency_vs_size:129:p99_us:1|.0f}\\
8 KiB & Native & r1 & \N{RWC:latency_vs_size:122:p50_us:1|.2f}/\N{RWC:latency_vs_size:122:p99_us:1|.0f} & \N{RWC:latency_vs_size:128:p50_us:1|.2f}/\N{RWC:latency_vs_size:128:p99_us:1|.0f} & \N{RWC:latency_vs_size:134:p50_us:1|.2f}/\N{RWC:latency_vs_size:134:p99_us:1|.0f}\\
8 KiB & Owner TCP & r0 & \N{RWC:latency_vs_size:118:p50_us:1|.0f}/\N{RWC:latency_vs_size:118:p99_us:1|.0f} & \N{RWC:latency_vs_size:124:p50_us:1|.0f}/\N{RWC:latency_vs_size:124:p99_us:1|.0f} & \N{RWC:latency_vs_size:130:p50_us:1|.0f}/\N{RWC:latency_vs_size:130:p99_us:1|.0f}\\
8 KiB & Owner TCP & r1 & \N{RWC:latency_vs_size:121:p50_us:1|.0f}/\N{RWC:latency_vs_size:121:p99_us:1|.0f} & \N{RWC:latency_vs_size:127:p50_us:1|.0f}/\N{RWC:latency_vs_size:127:p99_us:1|.0f} & \N{RWC:latency_vs_size:133:p50_us:1|.0f}/\N{RWC:latency_vs_size:133:p99_us:1|.0f}\\
8 KiB & Owner verbs & r0 & \N{RWC:latency_vs_size:119:p50_us:1|.0f}/\N{RWC:latency_vs_size:119:p99_us:1|.0f} & \N{RWC:latency_vs_size:125:p50_us:1|.0f}/\N{RWC:latency_vs_size:125:p99_us:1|.0f} & \N{RWC:latency_vs_size:131:p50_us:1|.0f}/\N{RWC:latency_vs_size:131:p99_us:1|.0f}\\
8 KiB & Owner verbs & r1 & \N{RWC:latency_vs_size:120:p50_us:1|.0f}/\N{RWC:latency_vs_size:120:p99_us:1|.0f} & \N{RWC:latency_vs_size:126:p50_us:1|.0f}/\N{RWC:latency_vs_size:126:p99_us:1|.0f} & \N{RWC:latency_vs_size:132:p50_us:1|.0f}/\N{RWC:latency_vs_size:132:p99_us:1|.0f}\\
64 KiB & Native & r0 & \N{RWC:latency_vs_size:135:p50_us:1|.0f}/\N{RWC:latency_vs_size:135:p99_us:1|.0f} & \N{RWC:latency_vs_size:141:p50_us:1|.0f}/\N{RWC:latency_vs_size:141:p99_us:1|.0f} & \N{RWC:latency_vs_size:6:p50_us:1|.0f}/\N{RWC:latency_vs_size:6:p99_us:1|.0f}\\
64 KiB & Native & r1 & \N{RWC:latency_vs_size:140:p50_us:1|.1f}/\N{RWC:latency_vs_size:140:p99_us:1|.0f} & \N{RWC:latency_vs_size:146:p50_us:1|.0f}/\N{RWC:latency_vs_size:146:p99_us:1|.0f} & \N{RWC:latency_vs_size:149:p50_us:1|.0f}/\N{RWC:latency_vs_size:149:p99_us:1|.0f}\\
64 KiB & Owner TCP & r0 & \N{RWC:latency_vs_size:136:p50_us:1|.0f}/\N{RWC:latency_vs_size:136:p99_us:1|.0f} & \N{RWC:latency_vs_size:142:p50_us:1|.0f}/\N{RWC:latency_vs_size:142:p99_us:1|.0f} & \N{RWC:latency_vs_size:7:p50_us:1|.0f}/\N{RWC:latency_vs_size:7:p99_us:1|.0f}\\
64 KiB & Owner TCP & r1 & \N{RWC:latency_vs_size:139:p50_us:1|.0f}/\N{RWC:latency_vs_size:139:p99_us:1|.0f} & \N{RWC:latency_vs_size:145:p50_us:1|.0f}/\N{RWC:latency_vs_size:145:p99_us:1|.0f} & \N{RWC:latency_vs_size:148:p50_us:1|.0f}/\N{RWC:latency_vs_size:148:p99_us:1|.0f}\\
64 KiB & Owner verbs & r0 & \N{RWC:latency_vs_size:137:p50_us:1|.0f}/\N{RWC:latency_vs_size:137:p99_us:1|.0f} & \N{RWC:latency_vs_size:143:p50_us:1|.0f}/\N{RWC:latency_vs_size:143:p99_us:1|.0f} & \N{RWC:latency_vs_size:8:p50_us:1|.0f}/\N{RWC:latency_vs_size:8:p99_us:1|.0f}\\
64 KiB & Owner verbs & r1 & \N{RWC:latency_vs_size:138:p50_us:1|.0f}/\N{RWC:latency_vs_size:138:p99_us:1|.0f} & \N{RWC:latency_vs_size:144:p50_us:1|.0f}/\N{RWC:latency_vs_size:144:p99_us:1|.0f} & \N{RWC:latency_vs_size:147:p50_us:1|.0f}/\N{RWC:latency_vs_size:147:p99_us:1|.0f}\\
\bottomrule\end{tabular}
\caption{Cold reads; successful API p50/p99 (\us{}), one row per repetition, never pooled. All three caller counts are shown. Values below 10 use two decimals, below 100 one, otherwise whole microseconds. Coordinator cache off. Unqualified cold cells have no latency. Unmatched dimensions: sync/grouping, RPC/security, index/routing, and host SSD (Native coordinator SSD versus owner SSD). Read lanes: Native 0, owner 1 at c1/c5 and 4 at c16; c5 owner reads are lane-limited, not a transport ranking. Floors are separate populations.}
\end{table}

\RWCFINDING{cold-populations}

\clearpage
\subsection{Warm reads and the cache contrast}
\begin{table}[H]\centering\small\setlength{\tabcolsep}{3pt}
\begin{tabular}{lrrrrr}\toprule
Size & Arm & Rep. & c1 & c5* & c16\\\midrule
1 KiB & Native & r0 & \N{RWC:latency_vs_size:3:p50_us:1|.2f}/\N{RWC:latency_vs_size:3:p99_us:1|.2f} & \N{RWC:latency_vs_size:153:p50_us:1|.2f}/\N{RWC:latency_vs_size:153:p99_us:1|.2f} & \N{RWC:latency_vs_size:159:p50_us:1|.2f}/\N{RWC:latency_vs_size:159:p99_us:1|.2f}\\
1 KiB & Native & r1 & \N{RWC:latency_vs_size:152:p50_us:1|.2f}/\N{RWC:latency_vs_size:152:p99_us:1|.2f} & \N{RWC:latency_vs_size:158:p50_us:1|.2f}/\N{RWC:latency_vs_size:158:p99_us:1|.2f} & \N{RWC:latency_vs_size:164:p50_us:1|.2f}/\N{RWC:latency_vs_size:164:p99_us:1|.2f}\\
1 KiB & Owner TCP & r0 & \N{RWC:latency_vs_size:5:p50_us:1|.1f}/\N{RWC:latency_vs_size:5:p99_us:1|.1f} & \N{RWC:latency_vs_size:155:p50_us:1|.0f}/\N{RWC:latency_vs_size:155:p99_us:1|.0f} & \N{RWC:latency_vs_size:161:p50_us:1|.1f}/\N{RWC:latency_vs_size:161:p99_us:1|.0f}\\
1 KiB & Owner TCP & r1 & \N{RWC:latency_vs_size:150:p50_us:1|.1f}/\N{RWC:latency_vs_size:150:p99_us:1|.1f} & \N{RWC:latency_vs_size:156:p50_us:1|.0f}/\N{RWC:latency_vs_size:156:p99_us:1|.0f} & \N{RWC:latency_vs_size:162:p50_us:1|.1f}/\N{RWC:latency_vs_size:162:p99_us:1|.0f}\\
1 KiB & Owner verbs & r0 & \N{RWC:latency_vs_size:4:p50_us:1|.1f}/\N{RWC:latency_vs_size:4:p99_us:1|.1f} & \N{RWC:latency_vs_size:154:p50_us:1|.0f}/\N{RWC:latency_vs_size:154:p99_us:1|.0f} & \N{RWC:latency_vs_size:160:p50_us:1|.1f}/\N{RWC:latency_vs_size:160:p99_us:1|.0f}\\
1 KiB & Owner verbs & r1 & \N{RWC:latency_vs_size:151:p50_us:1|.1f}/\N{RWC:latency_vs_size:151:p99_us:1|.1f} & \N{RWC:latency_vs_size:157:p50_us:1|.0f}/\N{RWC:latency_vs_size:157:p99_us:1|.0f} & \N{RWC:latency_vs_size:163:p50_us:1|.1f}/\N{RWC:latency_vs_size:163:p99_us:1|.0f}\\
8 KiB & Native & r0 & \N{RWC:latency_vs_size:165:p50_us:1|.2f}/\N{RWC:latency_vs_size:165:p99_us:1|.2f} & \N{RWC:latency_vs_size:171:p50_us:1|.2f}/\N{RWC:latency_vs_size:171:p99_us:1|.2f} & \N{RWC:latency_vs_size:177:p50_us:1|.2f}/\N{RWC:latency_vs_size:177:p99_us:1|.2f}\\
8 KiB & Native & r1 & \N{RWC:latency_vs_size:170:p50_us:1|.2f}/\N{RWC:latency_vs_size:170:p99_us:1|.2f} & \N{RWC:latency_vs_size:176:p50_us:1|.2f}/\N{RWC:latency_vs_size:176:p99_us:1|.2f} & \N{RWC:latency_vs_size:182:p50_us:1|.2f}/\N{RWC:latency_vs_size:182:p99_us:1|.2f}\\
8 KiB & Owner TCP & r0 & \N{RWC:latency_vs_size:167:p50_us:1|.1f}/\N{RWC:latency_vs_size:167:p99_us:1|.1f} & \N{RWC:latency_vs_size:173:p50_us:1|.0f}/\N{RWC:latency_vs_size:173:p99_us:1|.0f} & \N{RWC:latency_vs_size:179:p50_us:1|.0f}/\N{RWC:latency_vs_size:179:p99_us:1|.0f}\\
8 KiB & Owner TCP & r1 & \N{RWC:latency_vs_size:168:p50_us:1|.1f}/\N{RWC:latency_vs_size:168:p99_us:1|.1f} & \N{RWC:latency_vs_size:174:p50_us:1|.0f}/\N{RWC:latency_vs_size:174:p99_us:1|.0f} & \N{RWC:latency_vs_size:180:p50_us:1|.0f}/\N{RWC:latency_vs_size:180:p99_us:1|.0f}\\
8 KiB & Owner verbs & r0 & \N{RWC:latency_vs_size:166:p50_us:1|.1f}/\N{RWC:latency_vs_size:166:p99_us:1|.1f} & \N{RWC:latency_vs_size:172:p50_us:1|.0f}/\N{RWC:latency_vs_size:172:p99_us:1|.0f} & \N{RWC:latency_vs_size:178:p50_us:1|.1f}/\N{RWC:latency_vs_size:178:p99_us:1|.0f}\\
8 KiB & Owner verbs & r1 & \N{RWC:latency_vs_size:169:p50_us:1|.1f}/\N{RWC:latency_vs_size:169:p99_us:1|.1f} & \N{RWC:latency_vs_size:175:p50_us:1|.0f}/\N{RWC:latency_vs_size:175:p99_us:1|.0f} & \N{RWC:latency_vs_size:181:p50_us:1|.1f}/\N{RWC:latency_vs_size:181:p99_us:1|.0f}\\
64 KiB & Native & r0 & \N{RWC:latency_vs_size:183:p50_us:1|.1f}/\N{RWC:latency_vs_size:183:p99_us:1|.1f} & \N{RWC:latency_vs_size:189:p50_us:1|.1f}/\N{RWC:latency_vs_size:189:p99_us:1|.1f} & \N{RWC:latency_vs_size:9:p50_us:1|.1f}/\N{RWC:latency_vs_size:9:p99_us:1|.1f}\\
64 KiB & Native & r1 & \N{RWC:latency_vs_size:188:p50_us:1|.1f}/\N{RWC:latency_vs_size:188:p99_us:1|.1f} & \N{RWC:latency_vs_size:194:p50_us:1|.1f}/\N{RWC:latency_vs_size:194:p99_us:1|.1f} & \N{RWC:latency_vs_size:197:p50_us:1|.1f}/\N{RWC:latency_vs_size:197:p99_us:1|.1f}\\
64 KiB & Owner TCP & r0 & \N{RWC:latency_vs_size:185:p50_us:1|.1f}/\N{RWC:latency_vs_size:185:p99_us:1|.0f} & \N{RWC:latency_vs_size:191:p50_us:1|.0f}/\N{RWC:latency_vs_size:191:p99_us:1|.0f} & \N{RWC:latency_vs_size:11:p50_us:1|.0f}/\N{RWC:latency_vs_size:11:p99_us:1|.0f}\\
64 KiB & Owner TCP & r1 & \N{RWC:latency_vs_size:186:p50_us:1|.1f}/\N{RWC:latency_vs_size:186:p99_us:1|.0f} & \N{RWC:latency_vs_size:192:p50_us:1|.0f}/\N{RWC:latency_vs_size:192:p99_us:1|.0f} & \N{RWC:latency_vs_size:195:p50_us:1|.0f}/\N{RWC:latency_vs_size:195:p99_us:1|.0f}\\
64 KiB & Owner verbs & r0 & \N{RWC:latency_vs_size:184:p50_us:1|.1f}/\N{RWC:latency_vs_size:184:p99_us:1|.1f} & \N{RWC:latency_vs_size:190:p50_us:1|.0f}/\N{RWC:latency_vs_size:190:p99_us:1|.0f} & \N{RWC:latency_vs_size:10:p50_us:1|.0f}/\N{RWC:latency_vs_size:10:p99_us:1|.0f}\\
64 KiB & Owner verbs & r1 & \N{RWC:latency_vs_size:187:p50_us:1|.1f}/\N{RWC:latency_vs_size:187:p99_us:1|.1f} & \N{RWC:latency_vs_size:193:p50_us:1|.0f}/\N{RWC:latency_vs_size:193:p99_us:1|.0f} & \N{RWC:latency_vs_size:196:p50_us:1|.0f}/\N{RWC:latency_vs_size:196:p99_us:1|.0f}\\
\bottomrule\end{tabular}
\caption{Read preparation: Class B write settings, not a read durability guarantee. Warm reads; successful API p50/p99 (\us{}), repetitions separate; c1/c5/c16. Coordinator cache off. Unqualified cold cells have no latency. Unmatched dimensions: sync/grouping, RPC/security, index/routing, and host SSD (Native coordinator SSD versus owner SSD). Read lanes: Native 0, owner 1 at c1/c5 and 4 at c16; c5 owner reads are lane-limited, not a transport ranking. Floors are separate populations.}
\end{table}

\begin{figure}[H]\centering\RWCfigure{cold-warm}
\caption{Cold versus prewarmed whole-call means at one reader. Each point joins arm, size and repetition, but dataset geometry and operation counts differ. Shape identifies size; filled/open means r0/r1. No subtraction identifies SSD service.}
\end{figure}
\RWCFINDING{warm}

\clearpage
\subsection{Finite-case rates and the read-lane limit}
\begin{figure}[H]\centering\RWCfigure{throughput}
\caption{Successful API calls/s, not production capacity. c5 owner read points are shaded and lane-limited: one owner lane at c1/c5, four at c16; Native is in process. Do not rank transports on c5 cells. Write acknowledgement and cold/warm read contracts remain separate. Per-case in-flight, offered and unsuccessful counts remain in the CSV and figure evidence; missing occupancy is not zero.}
\end{figure}
\RWCFINDING{lanes}

\clearpage
\subsection{Floors, not an additive decomposition}
\begin{figure}[H]\centering\RWCfigure{floors}
\caption{Client/owner SSD and network floor means beside owner whole-call means, not a latency stack. Half RTT is doubled only where the receipt specifies it. Failed \texttt{ib\_write\_lat} event controls are unavailable, not zero; see Appendix~\ref{app:rwc-floors}.}
\end{figure}
\textbf{[INFERENCE]} Floors motivate hardware/software hypotheses but do not identify additive overhead. An exact partition needs disjoint, call-linked intervals from the same population; these independent controls do not supply them. No software remainder is computed.

Compare the \emph{same kind of call}: Native all-file fsync plus directory differs from owner one-file fdatasync. Append controls include fresh allocation; overwrite-fdatasync does not. Direct QD1 random reads differ from Native mapped/positional reads and owner exact-range buffered reads. Network tools do not include full RPC/security processing. Both host floors matter; equal VM class is not identical hardware.

\textbf{Limits.} One fleet and r0/r1 do not establish fleet spread or significance. Read/network floors have r0--r2; append controls r0/r1; three diagnostic profiles run once and remain separate. This run does not test five owners, NVMe-oF reads, Spot/other machine types, Sui replay, production capacity or failover. Three event-write floor failures are tool/mode failures, not system failures. The terminal report marks the campaign incomplete because of those failures; all planned cases have terminal receipts. No first-attempt result is used.
